\begin{proposition}[Deformations over a fixed genus-two curve]\hypertarget{tag-01GH}{}\label{cdp:st:E10}
\StatementRecord{cdp:st:E10}{proved}{01GH}{\SeriesRef{cdp:eq:E-Leray-map}, \SeriesRef{cdp:st:E5}}

The assertion that every small deformation of a fibration extends
over its original fixed positive-genus target curve $C$ is false even for a torus
fibration with Kodaira--Spencer rank two at a regular value.
\end{proposition}


\begin{proof}

\begin{enumerate}[label=\arabic*.,leftmargin=*]
\item Let $S\to\mathbb P^1_t$ be the rational elliptic surface obtained
by blowing up the nine base points of the Hesse pencil
$x^3+y^3+z^3-3txyz=0$; see \citet[\S2, pp.~238--241]{ArtebaniDolgachev2009}. Its critical values are
$\{1,\omega,\omega^2,\infty\}$, where $\omega^3=1$, $\omega\ne1$.
Pull another copy back by $\phi(t)=1/(t-4)$; its critical values
are $\{4,5,4+\omega,4+\omega^2\}$, disjoint from the first set.
Their fibre product $Z\to\mathbb P^1$ is a smooth projective
threefold, since at every point at least one elliptic-surface
map is a submersion. Its general fibre is a product of two
elliptic curves. Both elliptic $j$-functions are nonconstant,
so their two Kodaira--Spencer derivatives are simultaneously
nonzero off a finite set. The tensor of their product is then
diagonal of rank two.

\item Let $C_s$ be the smooth genus-two double cover
\[
 y^2=(t-10)(t-11)(t-12)(t-13)(t-14)(t-15-s)
\]
of $\mathbb P^1$, for $s$ in a small disc. All its branch values
avoid the critical values of $Z$. Thus
$Y_s=Z\times_{\mathbb P^1}C_s$ is a smooth projective deformation
of a compact threefold, and $g_s:Y_s\to C_s$ has connected torus
general fibre and rank two at a regular value. The Kodaira--Spencer
class $\theta_C\in H^1(C_0,T_{C_0})$ is nonzero: any infinitesimal
isomorphism of the hyperelliptic covers descends through their
unique hyperelliptic involutions to an infinitesimal projective
transformation fixing five distinct branch values; it is zero,
whereas the sixth branch value moves with derivative one.

\item Put
$L=\pi_C^*\mathcal O_{\mathbb P^1}(1)$ on $C_0$; the
double-cover canonical formula gives $K_{C_0}=L$, of degree two.
For either rational elliptic surface, relative duality gives
$R^0\pi_*\mathcal O=\mathcal O$ and
$R^1\pi_*\mathcal O=\mathcal O(-1)$: its relative canonical
line is $\pi^*\mathcal O(1)$ and every fibre has
$h^\bullet(\mathcal O)=(1,1)$. Relative K\"unneth and the flat
base change to $C_0$ consequently give
\[
 R^qg_*\mathcal O_Y=
 \begin{cases}
  \mathcal O_{C_0}&q=0,\\
  L^{-1}\oplus L^{-1}&q=1,\\
  L^{-2}&q=2,\\
  0&q>2.
 \end{cases}
\]
The vertical tangent sheaf is $g^*(L^{-1}\oplus L^{-1})$:
the invariant vector on each elliptic factor generates its
relative tangent sheaf, also at a node, where the kernel of
$d(xy)$ is generated by $x\partial_x-y\partial_y$.
Riemann--Roch on the genus-two curve gives
$h^0(L^{-a})=0$ and $h^1(L^{-a})=2a+1$ for $a>0$.
Projection and Leray give
\[
\begin{array}{c|rrrr}
 \mathcal F&h^0&h^1&h^2&h^3\\\hline
 \mathcal O_Y&1&2&6&5\\
 g^*T_{C_0}&0&3&10&7\\
 T_{Y/C_0}&0&6&20&14
\end{array}
\]
All higher groups vanish. In particular
$H^0(C_0,T_{C_0}\otimes R^1g_*\mathcal O_Y)=0$ and the
injection $H^1(C_0,T_{C_0})\to H^1(Y,g^*T_{C_0})$ is an
isomorphism between three-dimensional spaces.

\item Compatibility of the \v{C}ech transition equations for the family
$g_s$ gives
\[
 dg(\theta_Y)=g^*\theta_C
       \quad\text{in }H^1(Y_0,g^*T_{C_0}).
\]
The right side is nonzero because the first map of
\eqref{cdp:eq:E-Leray-map} is injective. Hence the first-order
obstruction to extending $g_0$ with fixed target $C_0$ is nonzero.
The obstruction belongs to
the base-deformation summand of \ref{cdp:st:E5}.

\end{enumerate}
\end{proof}